Magnetic actuation can steer a flexible catheter without placing motors along its body. A field applies torque to a magnetic segment, and a field gradient can supply force; elastic deformation and dissipation then determine the motion of the entire device \citep{abbott,kim}. For operation inside a confined passage, a useful dynamical model must therefore do more than predict the tip. It must connect available magnetic inputs to bounded centreline motion. This paper develops that connection for a clamped--free catheter model, using a learned Hamiltonian as the common object for numerical integration, nonlinear stability analysis and feedback design.

The application motivation is remote manipulation in reduced gravity. Spaceflight creates medical and operational conditions that warrant dedicated investigation; NASA's assessment of spaceflight-associated neuro-ocular syndrome is one example of a gravity-dependent health risk \citep{nasa}. For a remotely manipulated device in this setting, we isolate the following mechanical question: how to sustain a prescribed class of bounded motions and respect a clearance limit when a flexible magnetic device has distributed inertia, elasticity and damping. The model is planar, has small slopes and neglects gravity and contact. These assumptions make the energy balance explicit and permit independent finite-element calculations, while identifying the physical effects that a later device study must supply.

Two operating objectives organize the analysis. The first is recurrent motion with a proved amplitude bound, represented by an invariant torus of a conservative model. The second is tracking a requested endpoint path while keeping the intervening body inside a corridor. The first objective asks whether motion persists for arbitrarily long times; the second asks how much external work can safely be supplied during a manoeuvre. A common mechanical energy links them, but they call for different conclusions. We use KAM theory for the recurrent motion and an input-work inequality for clearance, then derive magnetic commands that make those conclusions relevant to the same damped beam.

The main difficulty is preserving this connection through learning and actuation. Low prediction error alone leaves the Hamiltonian's nonlinear derivatives undetermined. Even with accurate derivatives, a numerical solver can introduce an error comparable to the difference between two learned models. Finally, damping removes the recurrent motion of a conservative system unless a controller supplies the appropriate work. We address these issues in sequence: construct a power-consistent learned energy, resolve its approximation and integration errors, establish an invariant motion of the fitted function, and realize that motion with the available magnetic force and torque.

Recent magnetic-robot control studies address shape regulation and lumen navigation directly. \citet{francescon2025shape} combine a pseudo-rigid-link model with vision feedback and homogeneous fields to control a robot with a prescribed lengthwise magnetization profile, demonstrating navigation in a soft pancreatic-duct phantom under varied insertion poses. \citet{tong2025navigation} use a geometrically nonlinear magnetic-rod simulation with reduced contact handling to compute fields that reduce distal-tip constraint forces, with planar phantom experiments and three-dimensional numerical demonstrations. These studies establish practical routes from a mechanical model to magnetic commands. The question pursued here is how a learned energy, with a quantified identification error, can support both persistent motion and bounds on the entire centreline. For the stated planar beam, we connect analytic error and twist bounds to a finite-amplitude invariant torus, realize its low-mode dynamics in the damped full finite-element plant, and bound driven motion through supplied work. Table~\ref{tab:core-roadmap} summarizes the objects passed between these steps.

\begin{table*}[t]
\centering
\caption{Four steps of the study. Each step supplies the object needed by the next; the detailed proofs and additional experiments are in the Supplement.}
\label{tab:core-roadmap}
\begin{tabular}{p{0.16\textwidth}p{0.29\textwidth}p{0.44\textwidth}}
\toprule
Step & Question & Result passed to the next step\\
\midrule
1. Derive, Section~\ref{sec:derivation} & Which energy and discrete update describe the beam? & A power-consistent complex representation and a complete central symplectic update.\\
2. Evaluate, Section~\ref{sec:analysis} & What errors remain after fitting and integration? & Paired comparisons and quantitative approximation bounds for a specified learned function.\\
3. Establish persistence, Section~\ref{sec:kam} & Does that function preserve a nonlinear recurrent motion? & Nondegenerate twist, a finite-amplitude invariant torus and a local family of tori.\\
4. Realize, Section~\ref{sec:application} & How does the motion guide magnetic actuation? & Force/torque feedback, full-model amplitude bounds and clearance-constrained tracking.\\
\bottomrule
\end{tabular}
\end{table*}

\subsection{Learning a mechanical model that can be analyzed}
The starting point is a mechanical energy with explicit boundary work. Variational beam and rod formulations relate distributed strain energy, inertia and applied loads before discretization \citep{simo,antman}. In the present specialization, the tip magnet contributes a cosine potential under a static bias field. Its curvature enters the linear modal frequencies, while its higher-order terms govern the amplitude dependence of those frequencies. We retain the complete cosine because the nonlinear part later supplies the twist required for KAM nondegeneracy. Modal coordinates are normalized once, and the same dimensional map is used to interpret learned coefficients, simulated trajectories and magnetic commands. Supplement~\ref{S-sec:mechanics} supplies the distributed derivation; Section~\ref{sec:derivation} develops the model used in the analysis.

Hamiltonian neural networks provide a natural learning mechanism: fit a scalar energy and obtain its canonical vector field by differentiation \citep{greydanus}. This makes the energy derivatives available to both numerical integration and stability analysis. A magnetic catheter also loses and receives energy, so we follow the separation of conservative dynamics, dissipation and inputs used by dissipative SymODEN and port-Hamiltonian networks \citep{zhong,desai}. The learned conservative part specifies reversible dynamics; a nonnegative dissipation model and explicit magnetic ports complete the physical decoder. This separation is necessary for a controller to distinguish the energy stored by the device from the work it must supply.

Complex coordinates package modal position and momentum into one state. This representation supports shared complex features \citep{trabelsi}. Their physical interpretation remains tied to the underlying real energy. We derive that representation explicitly, including the conjugate dependence and the analytic extension used in the KAM argument. Complex-network approximation theory offers useful existence and comparison results, but those results depend on the activation, target class and norm \citep{zhang}. Here the relevant question is whether the fitted model controls the required derivatives on a declared domain. Smooth approximation theory \citep{hornik} provides an existence route; the quantitative analysis then uses the actual finite feature class and its fitted coefficients.

That feature class combines the complete magnetic potential with an analytic polynomial residual. Its coefficients can be fitted from energy derivatives or state increments, and all derivatives needed for the subsequent analysis are explicit. A residual may depend on momentum as well as position. We retain that dependence because it changes both the Hamiltonian vector field and the relation between physical velocity and canonical momentum. The magnetic feedback is consequently derived from the full learned gradient and Hessian. This gives learning a concrete role: it supplies a differentiable energy whose approximation error can be bounded and whose dynamics can be translated back into physical commands.

The distinction between identification and model reduction is particularly important for this device. A low-mode projection of a full beam trajectory need not satisfy a closed low-mode equation: the discarded modes also contribute to tip rotation and hence magnetic torque. A fitted residual can absorb part of that discrepancy on the observed trajectories. Its predictive value should therefore be assessed against independent full-beam solutions, while a theorem about that residual must use the function actually obtained by fitting. We organize the computations around both questions. Controlled reduced-model data isolate coefficient recovery and permit a quantitative learning-to-KAM argument. Independent full finite-element observations assess the discrepancy introduced by reduction and the response to driven motion.

The comparison design separates the effects of representation and mechanical knowledge. Real and complex implementations are fitted with and without the known quadratic energy, using paired training seeds and the same held-out trajectories. Five paired contrasts measure the effect of the prior in each representation, the representation effect with and without that prior, and their interaction. A separate integration comparison resolves numerical error before small learning differences are interpreted. The resulting real and complex predictions are equivalent at the resolved precision for this feature class.

\subsection{From numerical trajectories to persistent physical motion}
A learned Hamiltonian is continuous in time, whereas prediction and control require discrete updates. Central generating functions provide classical symplectic difference schemes for this step \citep{feng}. We derive the fourth-order central formula through its midpoint expansion and specialize it to the fitted magnetic energy. The correction, its symplecticity and its order belong to the established geometric-integration framework \citep{hairer}; the model-specific work is the complete realization for the learned, generally nonseparable Hamiltonian and its numerical evaluation. This treatment exposes the derivatives actually required by the update and allows the solver tolerance, phase error and state error to be assessed together.

The resulting numerical trajectories motivate the stability question, but KAM persistence is determined by properties of the Hamiltonian itself. Classical KAM theory connects sufficiently small perturbations of nondegenerate integrable systems to families of invariant tori \citep{poschel,bounemoura}. For the magnetic beam, the nonlinear frequency variation can be computed from the retained cosine. Learning must preserve that variation and control the full remainder. The analysis therefore follows coefficients and analytic norms through to the theorem assumptions.

Chen et al. connect nonzero Hamiltonian-network training loss to KAM theory under statistical and regularity hypotheses \citep{chen2022}. We develop a computable counterpart for the present analytic feature class. In the experiment used for the finite-amplitude proof, state-only COMSOL observations determine an integral regression. Bounds on its residual, integration defect and design matrix yield a coefficient-error bound; differentiating the same features converts it into analytic-error and twist bounds. Thus the learning error remains part of the argument even though the training loss is nonzero. The broader statistical route is retained in Supplement~\ref{S-sec:transfer}, while the main proof uses the deterministic bounds appropriate to these trajectory observations.

With those bounds in place, an a posteriori KAM theorem tests a computed torus of the fitted Hamiltonian through its invariance defect and geometric conditions \citep{haro2019aposteriori}. We apply it to the exported learned function, including its complete cosine and momentum-dependent residual. The result is a torus at approximately $1.9$ mm amplitude. Nondegeneracy of its final torsion also supplies a local family of positive measure through an established torus-stability result \citep{bounemoura2016torus}. Explicit measure estimates require a specified domain and constants \citep{biasco2018measure,chierchia2021measure}; we retain that calculation in the Supplement. Weighted Birkhoff averages additionally explore a wider domain over finite times \citep{das2017quasiperiodicity,duignan2023flows}, providing a complementary numerical view of the recurrent motions.

For these motions to guide actuation, the conservative subsystem must be connected to the damped full beam. We use a transverse magnetic force and torque to reproduce the learned low-mode accelerations. The remaining modes then form stable forced oscillators, whose quasiperiodic responses can be constructed and bounded. This yields an invariant torus of the complete $64$-mode model and a bound that also includes startup from zero higher-mode states. Leaving the torus for an endpoint task invokes a second control mechanism: a magnetic-work governor limits physical energy and hence the whole-centreline deflection. The two mechanisms address recurrent operation and manoeuvring within the same mechanical setting.

\subsection{Contributions and study design}
The work contributes three model-specific results. First, it derives the complex magnetic Hamiltonian, implements the classical fourth-order central update for the learned nonseparable energy, and resolves representation, prior and integration effects through matched comparisons. Second, it converts nonzero state-only identification error into computable analytic and twist bounds, then verifies a finite-amplitude invariant torus of that fitted function. Third, it constructs magnetic feedback and a full-model torus lift, together with an energy-based clearance condition for driven motion. The general approximation and KAM theorems are established results, cited with their hypotheses and identifiers where applied; the coefficient bounds, mechanical constructions and application proofs are developed here.

The quantitative results make the connection operational within the stated model. Starting the low state on the certified torus and all higher states at zero gives a whole-centreline bound below $1.971$ mm, a transverse-field bound below $0.922$ mT and a gradient bound below $0.02573$ T/m. A separate S-shaped endpoint task achieves $6.419\,\mu$m RMS error with the learned controller.

Two COMSOL identification settings serve distinct purposes within this common mechanical model: full finite-element observations test prediction and tracking, whereas the two-mode integral observations supply the learning-to-KAM proof and the associated magnetic realization. Their fitted coefficients and error budgets are specified separately. Section~\ref{sec:conclusion} brings the results together.
